% ============================================================
% Connected Minds Seed Grant Proposal
% (also suitable as a template for CIHR, NSERC, Brain Canada)
%
% Principal Investigator: Gunnar Blohm
% Trainee / Co-Investigator: Peter Ohue
% Institution: Queen's University, Centre for Neuroscience Studies
%
% TO EXPORT TO WORD: compile this .tex to PDF in Overleaf,
% then use Adobe Acrobat or https://www.adobe.com/acrobat/online/
% pdf-to-word.html to convert. Alternatively, open the .tex source
% in Overleaf and use the "Download Source" zip, then paste sections
% into Word manually.
% ============================================================

\documentclass[12pt]{article}
\usepackage[margin=2.5cm]{geometry}
\usepackage{setspace}
\usepackage{parskip}
\usepackage{titlesec}
\usepackage{booktabs}
\usepackage{tabularx}
\usepackage{array}
\usepackage{microtype}
\usepackage[numbers]{natbib}
\usepackage{hyperref}
\usepackage{enumitem}
\usepackage{graphicx}
\usepackage{amsmath}

\setlength{\parindent}{0pt}
\setlength{\parskip}{8pt}
\linespread{1.15}

\titleformat{\section}{\large\bfseries}{}{0em}{}[\titlerule]
\titleformat{\subsection}{\normalsize\bfseries}{}{0em}{}

\title{\textbf{Proposal for Connected Minds Seed Grant}\\[6pt]
\large\textbf{Neural Population Dynamics as a Bridge Between Computational
Motor Control and Adaptive Stroke Rehabilitation}}

\author{%
  Peter Ohue\textsuperscript{1} \and Gunnar Blohm\textsuperscript{1,2}\\[4pt]
  \small\textsuperscript{1}Centre for Neuroscience Studies, Queen's University,
  Kingston, ON, Canada\\
  \small\textsuperscript{2}Department of Biomedical and Molecular Sciences,
  Queen's University, Kingston, ON, Canada\\[4pt]
  \small Corresponding: \texttt{gunnar.blohm@queensu.ca}
}
\date{August 2026}

\begin{document}

\maketitle

% ============================================================
\section{Project Summary (250 words)}
% ============================================================

Motor recovery after stroke is among the most pressing unmet needs in
rehabilitation medicine.
Despite decades of research, fewer than half of stroke survivors recover
functional arm use, and existing therapies lack the mechanistic precision
needed to tailor training to the individual patient.
A central obstacle is that our ability to measure \emph{why} motor
adaptation succeeds or fails in a given session remains limited: conventional
clinical assessments capture task outcomes but not the internal neural
representations that determine whether a correction will be learned and
retained.

This proposal builds directly on our recent computational work demonstrating
that a reinforcement learning (RL) controller trained on adaptive
obstacle-avoidance develops structured, low-dimensional neural population
representations that encode perturbation context in a linearly accessible
form -- properties strikingly similar to those documented in biological
motor cortex during reaching.
We propose to extend this framework in two synergistic directions.

\textbf{Aim 1.} Characterise how neural population dynamics in the RL policy
network predict behavioural success and failure across perturbation
conditions, and test whether the representational geometry generalises to
novel tasks.

\textbf{Aim 2.} Translate the analytical pipeline into a closed-loop
virtual reality and wearable-sensor platform for upper-limb stroke
rehabilitation, using the RL agent's latent representations as real-time
biomarkers to personalise training difficulty and provide quantitative
outcome measures beyond conventional clinical scales.

The project directly advances the Connected Minds priority of
human-centred AI for health, generating tools and insights applicable to
any motor disorder in which adaptation under uncertainty is impaired.

% ============================================================
\section{Background and Rationale}
% ============================================================

\subsection{The Clinical Problem}

Stroke is a leading cause of persistent adult disability in Canada, with
approximately 62,000 new strokes occurring each year and upper-limb
motor impairment affecting the majority of survivors~\cite{HeartStroke2023}.
Recovery is highly heterogeneous: some patients regain near-normal function
while others plateau at severe impairment, and current clinical tools cannot
reliably predict which trajectory a given patient will follow
~\cite{Kwakkel2019}.
High-dose, task-oriented practice is the most evidence-supported
rehabilitation strategy~\cite{French2016}, but translating this principle
into individualized, adaptive protocols -- ones that automatically adjust
challenge level to the patient's changing capacity -- remains an open
engineering and scientific problem.

\subsection{The Scientific Gap}

The neural mechanisms underlying motor re-learning after stroke are poorly
understood at the population level.
Research in healthy adults has established that motor cortex activity during
reaching organises onto low-dimensional neural manifolds whose geometry
changes systematically with learning~\cite{Churchland2012,Sadtler2014,Oby2019}.
However, it is unknown whether similar manifold reorganisation occurs during
stroke recovery, whether the rate of that reorganisation predicts clinical
outcomes, or how external perturbation schedules should be designed to
accelerate it.

Our recent RL-based work provides a computational bridge.
We showed that a PPO controller trained purely by reward on an
obstacle-avoidance task spontaneously develops hidden-layer representations
with the same organisational properties -- low dimensionality, phase-aligned
clustering, graded kinematic encoding, linear decodability -- observed in
biological motor cortex~\cite{Ohue2026}.
Critically, the latent states of successfully correcting trials are
geometrically distinct from those of failure trials, suggesting that
population geometry is a principled indicator of adaptive capacity.
Translating this observation into a clinical tool requires: (i) extending
the RL framework to richer, human-relevant environments; and (ii) acquiring
a corresponding dataset from human stroke survivors performing analogous
reaching tasks.

\subsection{Connected Minds Alignment}

This proposal aligns with Connected Minds' focus on digital technologies
that augment human cognitive and physical capability.
The planned closed-loop VR-RL platform is an instance of human-centred AI:
the RL agent models the patient's evolving motor policy in real time, and
its latent representations drive an adaptive training protocol that is
responsive to the individual's neural and kinematic state rather than
following a fixed schedule.
The project also advances equity of access by targeting a platform that
is deployable in home settings, reducing dependence on specialized
rehabilitation centres.

% ============================================================
\section{Research Objectives and Hypotheses}
% ============================================================

\textbf{Overall objective:}
To establish low-dimensional latent neural population dynamics, derived
from both computational RL models and from kinematic/EMG surrogates in
human participants, as quantitative biomarkers of adaptive motor learning
and as drivers of individualized rehabilitation protocols.

\textbf{Hypothesis 1.}
The representational geometry of a trained RL policy's hidden-layer activity
will generalise across task variants: the same low-dimensional manifold
structure and phase-aligned clustering identified in the obstacle-avoidance
task will emerge in a richer perturbation environment with time-varying
force fields and dynamic obstacles.

\textbf{Hypothesis 2.}
In stroke survivors, the rate of manifold reorganisation during a
standardised perturbation-based reaching protocol -- quantified by the
change in population variance explained by the first two PCs across
sessions -- will correlate significantly with gains on the Fugl-Meyer
Upper Extremity (FMA-UE) scale at six-week follow-up.

\textbf{Hypothesis 3.}
A closed-loop training protocol in which task difficulty is modulated by
the RL agent's latent state geometry (specifically, by the Euclidean
distance of the current trial's latent trajectory from the learned
``success manifold'') will produce larger FMA-UE gains over six weeks than
a matched fixed-difficulty protocol, in a randomised single-site pilot trial.

% ============================================================
\section{Methodology}
% ============================================================

\subsection{Phase 1: Extended Computational Framework (Months 1--8)}

\textbf{Environment extension.}
We will extend the current two-dimensional point-mass environment to
include: (i) a three-degree-of-freedom planar arm with muscle-like
actuation and joint limits; (ii) time-varying force fields whose magnitude
and direction are drawn from the same distribution as the clinical
perturbation protocol; and (iii) dynamic obstacles that require online
re-planning.
The RL controller will be trained using PPO with the same hyperparameter
schedule validated in the baseline study, and the neural population
analysis pipeline (PCA, $K$-means clustering, linear decoding) will be
applied to the hidden-layer activity of the converged policy.

\textbf{Multi-seed validation.}
A minimum of ten independent training runs will be conducted to characterise
the variability in representational structure across seeds.
Features that are stable across seeds (low coefficient of variation) will
be designated as candidate biomarkers for translation to the clinical phase.

\textbf{Recurrent architecture comparison.}
A long short-term memory (LSTM) variant of the policy network will be
trained under identical conditions.
Comparing the manifold geometry between feedforward MLP and LSTM policies
will reveal whether temporal memory substantially changes the
representational organisation, informing the choice of architecture for
the closed-loop platform.

\subsection{Phase 2: Human Participant Study (Months 6--18)}

\textbf{Participants.}
We will recruit $n=20$ healthy adults and $n=30$ chronic stroke survivors
(FMA-UE 20--55, $\geq$6 months post-stroke, sufficient cognitive function
for VR interaction).
The sample size for the stroke group is powered to detect a correlation
of $r \geq 0.50$ between manifold reorganisation rate and FMA-UE change
($\alpha = 0.05$, $1-\beta = 0.80$, two-tailed).

\textbf{Instrumentation.}
Participants will perform a standardised perturbation-based reaching task
in a custom VR environment rendered on a head-mounted display (Meta Quest
Pro or equivalent).
Reaching kinematics will be captured at 100\,Hz using wrist-worn IMUs
(XSENS DOT) and processed to extract the same eight-dimensional state
vector used in the RL environment (goal-relative position, obstacle-relative
position, velocity).
Surface EMG will be recorded from six upper-limb muscles to provide a
neural surrogate for population analysis.

\textbf{Perturbation protocol.}
Seven perturbation conditions mirroring the computational study (P0, L1--L3,
R1--R3) will be applied via a lightweight cable-driven exoskeleton
(hand-exo or equivalent) or through VR-mediated haptic redirection.
Three sessions per week for six weeks will be administered, with each
session comprising 100 reaching trials.

\textbf{Population analysis.}
PCA and linear decoding will be applied to the EMG/kinematic population
vector at each timestep, paralleling the RL analysis pipeline.
The primary outcome for Hypothesis~2 is the session-by-session change in
$\%$ variance explained by PC1+PC2 of the EMG population, correlated
with FMA-UE change scores.

\subsection{Phase 3: Closed-Loop Pilot Trial (Months 14--24)}

\textbf{Adaptive protocol design.}
The difficulty-adaptation rule will use the latent state geometry from the
RL model as a reference manifold.
On each trial, the deviation of the participant's kinematic state trajectory
from the RL ``success manifold'' (projected into the common PCA space) will
be computed online.
If the deviation exceeds a threshold calibrated from the Phase 2 data, the
perturbation magnitude will be reduced by one step; if it falls below a
second threshold, the perturbation magnitude will be increased.
This rule mimics the bandwidth feedback used in motor learning research and
is consistent with the principles of optimal challenge point theory.

\textbf{Randomised pilot.}
$n=20$ stroke survivors will be randomised (1:1) to adaptive vs.
fixed-difficulty six-week protocols.
Primary outcome: FMA-UE change at six weeks.
Secondary outcomes: kinematic smoothness (SPARC), manifold reorganisation
rate, number of sessions required to reach each perturbation level,
patient-reported experience (NASA-TLX).

\textbf{Statistical analysis.}
The primary analysis is a mixed-model repeated-measures ANOVA with
group and session as factors.
Effect sizes (Cohen's $d$) will be reported with 95\% confidence intervals.
This is a feasibility and signal-detection trial; results will be used to
power a future definitive RCT.

% ============================================================
\section{Team and Expertise}
% ============================================================

\textbf{Gunnar Blohm (PI)} is a Professor and Canada Research Chair at
the Centre for Neuroscience Studies, Queen's University, with expertise
in computational motor neuroscience, sensorimotor transformations, and
neural population modelling.
His laboratory has a track record of developing and validating population-
level analyses of motor cortex and brainstem circuits.

\textbf{Peter Ohue (Graduate Trainee / Co-I)} is a doctoral candidate in the
Blohm laboratory with expertise in reinforcement learning, neural population
analysis, and computational neuroscience.
He is the primary developer of the RL-based obstacle-avoidance framework
on which this proposal is based, and will lead the computational and
experimental phases.

\textbf{Collaborators (to be confirmed):}
A rehabilitation medicine specialist and a certified occupational therapist
will provide clinical expertise for the stroke participant protocol and
outcome assessment.
A biomedical engineer with wearable-sensor expertise will support the IMU
and EMG data acquisition pipeline.

% ============================================================
\section{Expected Outcomes and Impact}
% ============================================================

\begin{enumerate}[leftmargin=*, label=\arabic*.]
  \item An extended RL framework with validated multi-seed representational
        stability, providing a robust computational foundation for clinical
        translation.
  \item A normative dataset of kinematic manifold dynamics in healthy adults
        and stroke survivors during standardised perturbation-based reaching,
        establishing reference values for the proposed biomarkers.
  \item Proof-of-concept evidence for the validity of RL latent geometry as
        a predictor of clinical rehabilitation outcomes.
  \item A pilot RCT comparing adaptive vs.\ fixed-difficulty VR-RL
        rehabilitation, with effect size estimates to power a future
        definitive trial.
  \item Open-source code and data deposited on Zenodo and GitHub, enabling
        independent replication and extension.
\end{enumerate}

The broader impact extends beyond stroke: the framework is applicable to
any motor disorder in which adaptation under perturbation is impaired,
including Parkinson's disease, cerebellar ataxia, and spinal cord injury.
The home-deployable VR-wearable platform addresses the access gap between
specialised rehabilitation centres and community care, directly advancing
Connected Minds' mandate of equitable, technology-enabled human
flourishing.

% ============================================================
\section{Timeline}
% ============================================================

\begin{tabularx}{\linewidth}{@{}lX@{}}
\toprule
\textbf{Period} & \textbf{Milestones} \\
\midrule
Months 1--4   & Environment extension (3-DOF arm, force fields, dynamic obstacles); multi-seed training; LSTM comparison \\
Months 3--6   & Ethics approval (human study); wearable instrumentation setup; pilot feasibility in healthy adults ($n=5$) \\
Months 6--12  & Healthy adult dataset ($n=20$); manifold analysis pipeline validation against RL reference \\
Months 10--18 & Stroke survivor observational study ($n=30$); Hypothesis~2 analysis \\
Months 14--24 & Adaptive-protocol pilot RCT ($n=20$ stroke survivors); Hypothesis~3 analysis \\
Months 22--24 & Manuscript preparation; data deposition; dissemination at FENS, SFN, and Canadian neuroscience meetings \\
\bottomrule
\end{tabularx}

% ============================================================
\section{Budget Justification}
% ============================================================

\textit{Note: Amounts are indicative for a Connected Minds Seed Grant
(typically \$50{,}000--\$100{,}000 over 1--2 years).
Adjust to the specific call ceiling and duration.}

\begin{tabularx}{\linewidth}{@{}lrX@{}}
\toprule
\textbf{Category} & \textbf{Amount (CAD)} & \textbf{Justification} \\
\midrule
Graduate student stipend (1 FTE, 24 months) & \$60{,}000 & Peter Ohue (doctoral trainee); standard Queen's University stipend (\$30{,}000/year) \\
Research assistant (part-time, 12 months)   & \$18{,}000 & Experimental support for human participant sessions \\
Wearable sensors (XSENS DOT $\times$4)      & \$8{,}000  & Kinematic capture for human participant protocol \\
VR hardware (Meta Quest Pro or equivalent)  & \$4{,}000  & Head-mounted display for VR reaching task \\
EMG system (Delsys Trigno $\times$6 ch.)    & \$6{,}000  & Neuromuscular signal acquisition \\
Computing (cloud GPU credits)               & \$3{,}000  & Multi-seed RL training; real-time latent state computation \\
Participant compensation (\$20/h, 50 h total) & \$4{,}000 & Healthy adults and stroke survivors \\
Open-access publication fees (2 papers)     & \$6{,}000  & PLOS Computational Biology or equivalent \\
Conference travel (2 conferences)           & \$4{,}000  & SFN, Canadian Association for Neuroscience \\
Indirect costs (20\%)                       & \$22{,}600 & University overhead rate \\
\midrule
\textbf{Total}                              & \textbf{\$135{,}600} & \\
\bottomrule
\end{tabularx}

\textit{For a \$50{,}000 ceiling: reduce to 12 months, one conference,
no EMG system; request the EMG as in-kind from the rehabilitation
collaborator.}

% ============================================================
\section{Data Management and Open Science}
% ============================================================

All experimental data will be de-identified and deposited on the Queen's
University institutional repository (QSPACE) and on Zenodo with a
persistent DOI within six months of collection.
Analysis code and trained RL models will be maintained on GitHub
(already public: \texttt{github.com/OhuePeter/NeuroRL-ObstacleAvoidance-v1.0})
under the MIT licence.
Manuscripts will be submitted to open-access journals or posted as
preprints on bioRxiv concurrent with journal submission.

% ============================================================
\section{Ethics and Risk}
% ============================================================

Human participant research will be conducted under Queen's University
Health Sciences and Affiliated Teaching Hospitals Research Ethics Board
(HSREB) approval, which will be sought in Months 3--5.
The perturbation protocol uses brief, low-force lateral impulses comparable
in magnitude to those used routinely in motor learning research and poses
negligible physical risk.
VR exposure will be limited to 30-minute sessions to minimise cybersickness.
Stroke participants will be screened to exclude severe spasticity, uncontrolled
epilepsy, and visual impairments.

% ============================================================
\section{Knowledge Translation}
% ============================================================

Short-term (\textless 2 years): peer-reviewed publications in computational
neuroscience and rehabilitation journals; open-source code and data release;
conference presentations.

Medium-term (2--5 years): submission of a full-scale CIHR Project Grant or
NSERC Discovery Grant building on the Seed results; regulatory-pathway
consultation for clinical device classification of the closed-loop platform.

Long-term ($>$5 years): commercialisation pathway in partnership with
Queen's Partnerships and Innovation (QPI) for a CE/FDA-clear home
rehabilitation device; integration with publicly funded rehabilitation
programmes in Ontario.

\bibliographystyle{unsrtnat}
\bibliography{references}

\end{document}
